\DTMsavedate{mydate}{2022-06-30}
\maketitle{Electronic Engineering - Circuits and Devices}{Electronic Engineering}{\DTMusedate{mydate}}

% ----------------------------------------------------- %
% COURSE INFO
\subsection*{Course Information}
\vspace{0ex}%
\noindent\parbox{0.5\textwidth}{%
    \noindent\begin{tabular}{@{}ll}
                 \textsf{Course:}          & Electronic Engineering - Circuits and Devices \\
                 \textsf{Course Code:}         & 	TEE 1203    \\
                 \textsf{Credits:}    & 	10
    \end{tabular}}
\vspace{2ex}\hrule\vspace{2ex}

% ----------------------------------------------------- %
% INSTRUCTOR INFO
\subsection*{Lecturer Information}
\noindent\parbox{0.5\textwidth}{%
    \noindent\begin{tabular}{@{}ll}
                 \textsf{Name:}         &    Svetlana A. Bebova           \\
% \textsf{Phone:}        &    \phone          \\
\textsf{Email:}        &    \href{mailto:swetlana.bebova@nust.ac.zw}{swetlana.bebova@nust.ac.zw} \\
\textsf{Qualifications:}        &  MSc in Radio Electronic Engineering (Technical University of Sofia)
    \end{tabular}}
\vspace{2ex}\hrule\vspace{2ex}


% ----------------------------------------------------- %
% COURSE SYNOPSIS
\subsection*{Course Synopsis}
The module outlines rectifying and Zener diodes: structure, operation, characteristics and parameters; Diode applications: rectifiers and power supplies; clippers and clamps; Schottky diode; Transistors; Bipolar Junction Transistors (BJTs): structure, terminals and operation; BJT configurations, static characteristics and parameters; biasing methods; d;c; circuit analysis and design; Darlington pair; BJT packages and data sheet; Field Effect Transistors (FETs): types, structure, terminals and operation; Configurations and static characteristics; d;c; circuit analysis; Power devices and heat sinks; Optoelectronic and photo- electronic devices: Light-Emitting Diodes (LEDs), infra-red diodes, 7-segment displays, Liquid Crystal Displays (LCDs), photodiode and phototransistor; Applications; Thermistors: structure, types, operation and applications.
\vspace{2ex} \hrule\vspace{2ex}
% ----------------------------------------------------- %
% COURSE TOPICS
\subsection*{Topics}
\begin{outline}
    \1 Semiconductor Diodes
    \1 Semiconductor Materials and BJT
    \1 Diodes Application
    \1 Bipolar Junction Transistor
    \1 BJT Biasing and equivalent circuits
    \1 Field Effect Transistors
\end{outline}
\vspace{2ex} \hrule\vspace{2ex}

% ----------------------------------------------------- %
% Students Assenssment
\subsection*{Assessment of Students}
\begin{outline}
    \1 Students are assessed through written assignments and written in-class tests that shall form part of course work assessment mark.
    \1 There shall be a final examination at the end of the semester.
\end{outline}
\vspace{2ex}\hrule\vspace{2ex}

% ----------------------------------------------------- %
% Grade Evaluation
\subsection*{Course Evaluation and Grading Policies}
Grades are based on:
Continuous assessment (\qty{25}{\percent}),
Final Exam (\qty{75}{\percent})
\vspace{2ex}\hrule\vspace{2ex}

% Policies and Other information
\subsection*{Policies and Other Information}
% Conduct
\subsection*{Key Text}
\begin{itemize}
    \item Electronic Principles, 8th Edition, 2015 by A. Malvino, D. Bates
    \item The Art of Electronics, 3rd Edition, 2015 by P. Horowitz
\end{itemize}
\vspace{2ex}\hrule
\vspace{2ex}\hrule\vspace{2ex}

